\subsection{相对可分代数闭包}

命题 12.——设 E 是 K 的一个扩张，\(E_{s}\) 为 E 中在 K 上代数且可分的元素所成之集。则 \(E_{s}\) 是 E 的在 K 上代数且可分的最大子扩张。

由命题 6，a)（V，第 39 页），E 的每个在 K 上代数且可分的子扩张都含于 \(E_{s}\) 。由命题 6，b)（同处），K 的扩张 \(\mathrm{K}(E_{s})\) 是代数且可分的，因此 \(\mathrm{K}(E_{s}) \subset E_{s}\) ，从而 \(\mathrm{K}(E_{s}) = E_{s}\) 。

沿用前一命题的记号，\(E_{s}\) 称为 K 在 E 中的相对可分（代数）闭包。当 K 为完全域时，\(E_{s}\) 是 K 在 E 中的相对代数闭包（V，第 36 页，命题 2）。

命题 13.——设 E 是 K 的一个代数扩张，并设 \(E_{s}\) 为 K 在 E 中的相对可分代数闭包。

\begin{enumerate}
\def\labelenumi{\alph{enumi})}
\item
  E 是一个 \(p\) -根式扩张（在 \(\mathrm{E}_s\) 上）。
\item
  若 \(\mathbf{F}\) 是 \(\mathbf{E}\) 的一个子扩张，使得 \(\mathbf{E}\) 是 \(p\) -根式的（在 \(\mathbf{F}\) 上），则 \(\mathbf{F} \supseteq \mathbf{E}_s\) .
\item
  \(\mathbf{E}_s\) 是 \(\mathbf{E}\) 的唯一在 \(\mathbf{K}\) 上可分且使 \(\mathbf{E}\) 在其上是 \(p\) -根式的子扩张。
\end{enumerate}

只需在下列情形证明 \(a\) )：\(\mathbf{K}\) 的特征为 \(p \neq 0\) 。设 \(x\) 为 \(\mathbf{E}\) 的元素，\(f\) 为其在 \(\mathbf{K}\) 上的极小多项式。存在整数 \(m \geqslant 0\) 使得 \(f\) 属于 \(\mathbf{K}[\mathbf{X}^{p^m}]\) 但不属于 \(\mathbf{K}[\mathbf{X}^{p^{m+1}}]\) ；换言之，我们有 \(f(\mathbf{X}) = g(\mathbf{X}^{p^m})\) ，其中 \(g \in \mathbf{K}[\mathbf{X}], g \notin \mathbf{K}[\mathbf{X}^p]\) 。由于 \(f\) 不可约，\(g\) 亦不可约，因此 \(g\) 是 \(x^{p^m}\) 在 \(\mathbf{K}\) 上的极小多项式。于是由 \(V\) ，第 38 页命题 4 及第 39 页命题 5，我们有 \(x^{p^m} \in \mathrm{E}_s\) ，故 \(\mathbf{E}\) 是 \(p\) -根式的（在 \(\mathrm{E}_s\) 上）.

现在假设 \(b)\) 成立，并设 \(x \in \mathrm{E}_s\) 。由于 \(x\) 在 \(\mathbf{K}\) 上可分，它在 \(\mathrm{F}(\mathrm{V}, \mathrm{p.~39}, \mathrm{Cor.~2})\) 上也可分；又由于 \(\mathrm{E}\) 是 \(p\) -根式的（在 \(\mathrm{F}\) 上），\(x\) 也是 \(p\) -根式的（在 \(\mathrm{F}\) 上），因此 \(x \in \mathrm{F}(\mathrm{V}, \mathrm{p.~39}, \mathrm{Cor.~3})\) .

最后，c) 由 a)、b) 及命题 12 得出。

推论 1.——设 E 与 K' 是包含在 K 的同一扩张中的 K 的两个扩张。假设 E 在 K 上代数，并记 \(E_{s}\) 为 K 在 E 中的相对可分代数闭包。则 \(\mathrm{K}'(\mathrm{E}_{s})\) 是 \(K'\) 在 \(\mathrm{K}'(\mathrm{E})\) 中的相对可分代数闭包。

因为 \(\mathbf{K}^{\prime}(\mathbf{E}_{s})\) 是 \(K^{\prime}\) 的可分代数扩张（命题 10，V，第 42 页）；由于 E 在 \(E_{s}\) 上是 p-根式的，\(\mathbf{K}^{\prime}(\mathbf{E})\) 在 \(\mathbf{K}^{\prime}(\mathbf{E}_{s})\) 上的扩张是 p-根式的（V，第 25 页，推论）。于是只需应用命题 13。

推论 2.——若 E 在 K 上次数有限，则 \(E_{s} = \bigcap_{n \geqslant 0} K(E^{p^{n}})\) .

对每个整数 \(n \geqslant 0\) ，用 \(F_n\) 记子扩张 \(K(E^{p^n})\) （\(E\) 的）。序列 \((F_n)_{n \geqslant 0}\) （为 \(E\) 的向量子空间）是递降的，且 \(E\) 在 \(K\) 上维数有限。因此存在整数 \(m \geqslant 0\) 使得 \(F_m = F_n\) （对一切 \(n \geqslant m\) ）。于是我们有 \(K(F_m^p) = F_{m+1} = F_m\) ，因此 \(F_m\) 是一个可分扩张 \(K(V, p. 42, \text{Prop. } 11)\) ；显然 \(E\) 是 \(p\) -根式的（在 \(F_m\) 上），故命题 13 蕴含 \(E_s = F_m = \bigcap_{n \geqslant 0} F_n\) .

注记.——设 E 是 K 的一个代数扩张，\(E_{r}\) 是 E 在 K 中的相对 p-根式闭包（V，第 25 页）。则 \(E_{r}\) 是 E 的在 K 上 p-根式的最大子扩张（V，第 25 页，命题 2）。然而，一般而言 E 在 \(E_{r}\) 上并不可分（V，第 152 页，习题 2）；关于拟伽罗瓦扩张的情形，见 V，第 76 页。

